\chapter{Research Methodology\label{chapter:methodology}}

The previous chapter set out what the study must achieve. This one sets out how: the research
strategy, the method chosen for each research question and why it fits, the controls that make a
language comparison fair, the statistical treatment of the measurements, and the reproducibility
and ethical commitments the work is carried out under.

\section{Research approach}\label{sec:approach}
The study is experimental with a constructive (design-science) element. The artifact is a pair
of functionally equivalent implementations of the subsystem, one in C++ and one in Rust, plus
the framework that benchmarks them, both built from a single frozen specification and evaluated
under controlled experiments. Evidence is triangulated across the three independent methods
Table~\ref{tab:design} lists, one of which, the FFI migration case study answering RQ4, is part
measured and part qualitative. A weakness in one of them, such as the
environment sensitivity of tail-latency measurement, therefore does not carry into the
conclusions as a whole. Every measurement reports variance over repeated runs, and
Section~\ref{sec:controls} sets out the controls that keep the language the only variable.

That design follows from the shape of the questions. Each asks how one outcome (tail latency,
defects prevented, resource use, migration cost) depends on one variable, the implementation
language, which therefore has to be the only systematic difference between the two objects under
test. That rules out the commonest language comparison in the literature, where two
independently written programs are measured and a gap carrying the design, the libraries and the
skill of their authors is attributed to the language. Building both sides from one specification
removes the confound: the comparison is between two encodings of one design, and the artifact is
both the contribution and the instrument that measures it. Where a question can be answered with
numbers the treatment is quantitative, and where it cannot, as with the kind of guarantee in RQ2
and the API-design effort in RQ4, it is qualitative but still systematic.

The gap this design addresses (Section~\ref{sec:gap}) is only as good as the search that looked
for it, so Appendix~\ref{appendix:search} records the databases, the term combinations, the
citation chaining and the filtering rules. Two of its findings bear on the design and not just
on the bibliography. The vendor-hosted theses closest to this study
(Section~\ref{sec:relatedwork}) appear in no citation database, so chaining never reaches them
and they had to be found in the institutional repositories directly. The same searches establish
an absence on the other side: no Finnish thesis on Rust in a radio-access or telecom context, and
no indexed publication mentioning Rust from the vendor hosting this work. The domain-specific
literature on the language question is therefore Swedish and Ericsson-hosted, whether it carries
over here is untested, and being grey literature it is cited as evidence of what a vendor has
attempted and found rather than as peer-reviewed results. No study, indexed or grey, measures
the tail-latency behaviour of a Rust and a C++ implementation of the same stateful, concurrent
control-plane subsystem under overload.

All four chains in Table~\ref{tab:design} end in the same place: Section~\ref{sec:rqdiscussion}
judges each question and Section~\ref{sec:hypotheses} disposes of its hypothesis. RQ1 and RQ3
share H1,
worded to cover performance equivalence in both the latency and the resource sense, and that
wording turns out to be the hypothesis's main weakness. RQ2 is insensitive to the environment,
and RQ4 is so except along one axis: build time is a timing measurement, read as a ratio between
two toolchains measured alternately on one machine. RQ1 and RQ3 depend on the measurement host
and are reported against what it can carry (Section~\ref{sec:testbed}).

\begin{table}[H]
\centering
\caption[Research design: how each research question is answered]{Research design: how each research question is answered. Every question is tied to a
concrete method, a published artifact, the data it produces, and the metric on which the
conclusion rests. RQ1 and RQ3 carry two methods rather than one: the sweep measures the
round-trip a deployment would experience, and the component decomposition measures the part of
it the language is responsible for. Section~\ref{sec:methodperf} states why one cannot stand in
for the other. RQ1 and RQ3 are tested against H1, RQ2 against H2, and RQ4 against
H3.\label{tab:design}}
\renewcommand{\arraystretch}{1.25}
\begin{tabular}{@{}p{0.06\linewidth} p{0.24\linewidth} p{0.24\linewidth} p{0.32\linewidth}@{}}
\hline\hline
RQ & Method & Artifact / data & Primary metric \\
\hline
RQ1 & Load sweep to overload, interleaved repeats & Benchmark harness, latency histograms &
median, p99/p99.9/p99.99 tail latency, and maximum vs.\ offered load \\
RQ3 & Same sweep, resource capture & \texttt{getrusage} samples per run &
peak resident set size (RSS), CPU utilisation, throughput \\
RQ1, RQ3 & Component decomposition: the core called directly, then through the service &
Matched microbenchmark pair, interleaved repeats &
per-operation cost split into the operation and the cross-thread handoff \\
RQ2 & Controlled fault injection & Defect catalogue, ASan/TSan/UBSan, \texttt{rustc},
Loom/Shuttle and Miri logs &
defect-class taxonomy (compile-time / runtime / undetected) \\
RQ4 & FFI migration prototype & C++ host $+$ Rust C-ABI seam &
seam size, \texttt{unsafe} scope, per-call interop cost \\
\hline
\end{tabular}
\end{table}

\section{Methods for the four research questions}

\subsection{Performance measurement under a load sweep (RQ1, RQ3)}\label{sec:methodperf}
RQ1 and RQ3 are answered by driving both implementations with one deterministic synthetic
workload across a sweep of offered load, from light load through saturation into overload. Two
design decisions make the measurement sound. The load generator follows an \emph{open-loop}
schedule: each request has a fixed intended send time, and latency is measured against that time
and not against the moment of admission, which corrects for coordinated omission as
Chapter~\ref{chapter:background} defines it. The sweep is pushed past saturation deliberately,
because the overload cliff appears only once offered load exceeds service capacity
(Section~\ref{sec:predictability}). Latency goes into high-dynamic-range histograms configured
identically on both sides, and resource use is sampled through the operating system.

\paragraph{Measuring the component as well as the round-trip.}
A round-trip bundles two costs of different character. One is the operation the language runs
inside a single thread, a hash lookup or a state-machine guard, all compiler-generated. The
other is the cross-thread handoff around it, an enqueue, a futex wake, a context switch, none of
it language-specific. Only the first can plausibly differ between C++ and Rust, and no
end-to-end measurement separates them. The method therefore works at two levels: the operation
processor called directly with no queue in the path, then the full threaded service, whose
difference is the handoff. Both are batch means over a long loop and not per-call medians,
because a clock read costs tens of nanoseconds,
the same order as the operation itself, so per-call timestamps would measure the instrument. The
languages are interleaved across repeats and compared with the same exact paired test as the
sweep, so the levels differ in what they measure and not in how they are analysed.

Reporting both separates two questions the research questions did not distinguish as originally
worded: What a deployment of this shape would experience, and what the language is doing.
Section~\ref{sec:budget} finds the two more than two orders of magnitude apart, so the
distinction decides whether the result can be interpreted at all. The second level was added
after a first pass of the campaign, and Section~\ref{sec:threats} records that revision.

\paragraph{Isolating the contended synchronisation primitives.}
The overload comparison in Section~\ref{sec:overload} attributes its gap to two forced-idiom
entries where the standard libraries diverge, the bounded blocking queue and the one-shot
completion handle. That attribution is argued from the shape of the full-service measurement, so
a third matched microbenchmark drives those two primitives alone, with no \texttt{ShardCore}
work in the path, under the same interleaved repeats and paired test.
Section~\ref{sec:contention} reports whether the isolated primitives reproduce the ordering the
full service showed.

\subsection{Controlled fault injection (RQ2)}\label{sec:methodfaults}
RQ2 asks which defect classes a language prevents, and measuring defect-free code cannot answer
that, because the absence of a bug in a passing test says nothing about whether the language
would have caught it. The method is therefore adversarial, following the established practice of
software fault injection, where faults are introduced deliberately to see how a system behaves
in their presence~\cite{natella2016fault}. A catalogue of representative C++ defects is written
\emph{deliberately} from the classes that affect this subsystem (Chapter~\ref{chapter:context}),
each one a small self-contained program flavoured with the domain, built and run under
AddressSanitizer, ThreadSanitizer and UndefinedBehaviorSanitizer, recording which tool detects
it, if any, and at what stage.

The same construct is then written in safe Rust and compiled. The compiler's outcome is
recorded, and where it compiles, so is the program's behaviour. The boundary classes give the
compiler nothing to say, so where another tool can reach one, the Rust side goes to that tool as
well: the two concurrency classes to the model checkers Loom and Shuttle, and the class inside
an \texttt{unsafe} block to the Miri interpreter. Those three are the cases these tools were run
on, chosen because each is the kind of defect the tool is built for, and answering ``at what
stage and at what cost'' for them is what the extra runs buy. Each is paired with a control in
which the defect is fixed, so that finding the defect and reporting everything are not the same
observation. Whether ThreadSanitizer reports the lock-order inversion depends on the schedule,
so the runner executes that case twelve times and reports the rate, and intermittent detection
is not mistaken for reliable detection.

\paragraph{Why the C++ side is tooled with sanitisers only.}
No static analyser is run against the C++ catalogue, and that is a decision. RQ2 does not ask
which tool finds the most defects. It asks at what \emph{stage}, and with what force, each
language rules a defect out, and the taxonomy's three outcomes (compile-time rejection, runtime
detection, silent miss) are the categories that question needs. An opt-in pass belongs to none
of them: a warning a build can be configured to ignore is a different kind of object from a
compilation that cannot succeed, and that difference is what the method measures. The cost is a
C++ side narrower than the best a C++ project could manage, since a well-configured analyser
could plausibly reach the small intra-procedural memory-safety cases, though not the concurrency
ones. The study therefore reports the tooling it ran, the sanitiser-based practice of
Section~\ref{sec:cpplatency}, and not a ceiling on what C++ tooling can achieve.

The injected fault is held fixed while the language varies, so any difference in outcome comes
from the language's guarantees and not from the specific bug, and what the method establishes is
the \emph{kind} of protection each language offers, not a performance number. One further
decision determines whether the taxonomy can be trusted: a catalogue drawn only from classes the
language is known to prevent could confirm the hypothesis but never disturb it, and would
describe the author's choice of cases as much as the language. It is therefore built in two
groups, classes inside Rust's stated guarantee and classes just outside it, where ownership and
the \texttt{Send}/\texttt{Sync} discipline give no protection, each case carrying a recorded
expectation the runner checks in either direction, so an unexpected outcome surfaces as a
discrepancy instead of being absorbed into the narrative. What this buys is a boundary and not a
score: which classes the language removes, and which comparable ones it leaves in place, which
is what a migration decision needs.

\subsection{Foreign-function-interface migration case study (RQ4)}\label{sec:methodffi}
RQ4 concerns the practical cost of adopting Rust inside an existing C++ system, an engineering
question as much as a performance one. It is answered with a prototype built in the realistic
direction of migration, a C++ host calling a Rust component across a C application binary
interface, on which three quantities are measured. The \emph{size} of the seam counts its entry
points, lines of code and dependencies. The \emph{scope} of \texttt{unsafe} records where the
checked language is left and whether the safety-critical core is affected. The \emph{per-call
cost} is the price of crossing the boundary, measured against the same operation invoked
natively so the comparison isolates what the seam adds. The first two are read off the code and
the third is a matched pair of microbenchmarks, and pairing them separates the run-time cost of
interoperability from its true cost, which lies in API design and human learning.

Interoperability is one of the three terms RQ4 names, and a prototype answers it alone. Three
further quantities cover the other two, chosen because a build engineer reaches them before
anything the seam can show. \emph{Build time} is measured clean and incrementally, for the core
library alone and for every default target, with the same job count and interleaved repeats on
both sides so that one host's load falls on both toolchains equally. \emph{Static footprint} is
measured on the conformance runner, the one executable both languages build from the same
specification, and at both linkage policies rather than at each toolchain's default, because the
defaults differ in whether the standard library is linked in and a single pair of numbers would
report that choice as a language difference. \emph{Third-party surface} counts the transitive
crates each configuration resolves and the \texttt{unsafe} sites inside them, which makes the
dependency-free claim of Section~\ref{sec:idioms} checkable. The first is host-sensitive and the
other two are not.

\section{Fairness as a controlled-variable strategy}\label{sec:controls}
Every research question attributes a difference to the language, so validity depends on holding
everything else equal, and fairness is therefore the study's central control, enforced at
several levels, which Table~\ref{tab:controls} states variable by variable. The
\emph{specification} is frozen before any code is written, down to the sharding and concurrency
model (Section~\ref{sec:fairness}), the \emph{run-time substrate} is equalised by allocator and
optimisation level as closely as the toolchains allow, and one deterministic \emph{workload}
driver puts both implementations through the same sequence of operations for a given seed and
active-UE population, so any latency difference reflects how each language executes the same
work and not a difference in what work is done. Where the two languages genuinely cannot express
the same construct, the deviation is recorded in a forced-idiom log. What remains uncontrolled,
principally single-author skill asymmetry, is acknowledged as a threat to validity
(Chapter~\ref{chapter:discussion}) and not claimed to be eliminated.

The table's third column matters most: a controlled experiment is credible in proportion to how
honestly it accounts for what it failed to control. Three rows record a residual and not a clean fix, since the languages cannot
always express the same construct, they do not ship the same hash-map layout, and one person
wrote both implementations. The forced-idiom log handles the first, and the last row names the
one variable with no control at all.

\begin{table}[H]
\centering
\small
\caption[The controlled-variable strategy]{The controlled-variable strategy: what is held fixed so that the language remains the
only systematic difference, and what each control does not reach. The third column is not a
list of oversights but the study's own account of its limits. Each entry is picked up again as
a threat to validity in Section~\ref{sec:threats}. Only the last row names a variable for
which there is no control at all.\label{tab:controls}}
\footnotesize
\renewcommand{\arraystretch}{1.15}
\begin{tabular}{@{}R{0.16\linewidth} R{0.44\linewidth} R{0.32\linewidth}@{}}
\hline\hline
Variable & How it is held fixed & What the control does not reach \\
\hline
Design & One specification, frozen before any code is written, fixing the public API, the state machine, the data structures and their algorithmic strategy, and how state is divided into shards (\S\ref{sec:fairness}) & Where the two languages cannot express the same construct, the divergence is logged, not removed (Table~\ref{tab:idioms}) \\
Concurrency model & Thread-per-shard over bounded blocking multi-producer, single-consumer queues on both sides, with no asynchronous runtime on either & Nothing outstanding \\
Hashing & SplitMix64 for shard selection and FNV-1a, the Fowler--Noll--Vo hash, inside the shard, pinned identically, so neither standard library's default hasher becomes an accidental advantage & The container itself is still each language's own hash map, and the two differ in layout (\S\ref{sec:budget}) \\
Allocator & The same \texttt{jemalloc} linked on both sides, with a system-allocator build kept on each as a control & Nothing outstanding \\
Optimisation & \texttt{-O2} against \texttt{opt-level~=~2} & Nominal parity is not identical code generation \\
Workload & One deterministic driver. For a given seed and active-UE population both implementations execute the same operation sequence, checked by per-kind counts & Synthetic and spec-justified rather than replayed from production traffic \\
Instrumentation & The same HDR histogram algorithm and configuration, the same open-loop schedule, the same repeat and interleaving discipline & Nothing outstanding \\
Behaviour & A language-neutral conformance suite whose canonical transcripts are diffed op for op (\S\ref{sec:parity}) & Parity bounds behaviour, not performance \\
\hline
Author & No control & One person wrote both implementations and is not equally expert in both \\
\hline
\end{tabular}
\end{table}

\section{Statistical treatment of measurements}\label{sec:stats}
Performance measurements are non-deterministic, so single-run numbers are avoided throughout.
Each configuration (language $\times$ allocator $\times$ load level) is measured over many
independent repeats, and the four variants are \emph{interleaved} within each repeat so that
slow drift, thermal effects, and background activity perturb them equally instead of biasing
one. Results are summarised by the median with a nonparametric bootstrap $95\%$ confidence
interval~\cite{efron1993bootstrap} rather than by the mean with a standard deviation, because
latency distributions are heavy-tailed and skewed and a distribution-free interval on the median
makes no normality assumption. That follows established guidance on statistically rigorous
performance evaluation, which warns that a single number invites conclusions run-to-run
variation does not support~\cite{georges2007rigorous}. Tail percentiles (p99, p99.9, p99.99) are
read from the HDR histograms, whose logarithmic bucketing retains three significant figures of
resolution across the whole range so that high percentiles are not quantised away. Intervals are
reported in full rather than as a symmetric half-width, because a bootstrap interval on the
median of a heavy-tailed sample is markedly asymmetric and a half-width wide enough to cover its
upper bound would place the lower bound below zero.

How far into the tail a campaign can see is a property of how many operations it recorded. At
the representative load a single run retires on the order of $170$~thousand measured
operations, so its p99.99 rests on roughly seventeen observations and is dominated by sampling
noise even on a perfectly quiet host. Reading the far tail therefore requires persisting each
run's complete HDR bucket contents, so that the repeats can be \emph{pooled} into one
distribution of several million operations. The harness does that, and keeps pooling and
averaging apart, because they are different statistics. The pooled figure takes the operation as
the unit of analysis and resolves the far tail, at the cost of carrying no across-repeat
interval, while the median-of-repeats figures reported with intervals keep the repeat as the
unit and remain the only basis for a paired comparison. The reported campaign's p99.99 is pooled
by that procedure and reaches Table~\ref{tab:latency} as a column of its own, marked because it
is the one column of that table computed at a different unit of analysis and so the one carrying
no interval. The analysis scripts read the pooled file and never the per-run field, so a
campaign whose histogram sidecars have not been pooled loses the column rather than printing the
per-run figure just called uninterpretable. Jitter is likewise measured directly rather than
inferred from the latency distribution, as the absolute difference between consecutive latencies
within a producer. That is the packet-delay-variation quantity of Request for Comments
(RFC)~3393~\cite{rfc3393}, and the sense in which this thesis uses the term: variation from one
operation to the next, not spread of the distribution as a whole. It is recorded and reported
descriptively, because no research question asks for it.

Neither p99.99 nor jitter enters the family of paired tests described below, and that follows
the same rule as the rest of the design rather than any shortage of data. Both are computed from
precisely the operations that interference lands on, which is the class of statistic pairing
does not rescue, so a paired interval on them could not be read as a language effect. They
belong in a descriptive column, and the family stays the size the design committed to.

Interval estimates alone are a weak basis for a comparison. Two intervals can overlap while one
variant is ahead in every single repeat, and they can separate on a difference that does not
reproduce. The study therefore tests the comparison directly, which the interleaved design makes
possible: because the four variants are measured back to back inside each repeat, the runs form
a matched design in which the interference common to a block cancels in the pairwise difference.
Language is compared against language with the allocator held fixed, using an exact two-sided
paired permutation test on the repeat-wise differences. The null hypothesis is that a pair
differs only by run-to-run noise, under which the sign of each paired difference is
exchangeable. For the repeat counts used here every one of the $2^{n}$ sign assignments is
enumerated rather than sampled, so the $p$-value is exact and, like the bootstrap, assumes
nothing about the shape of the distribution. Because the same comparison is made for several
metrics and both allocators, the family is corrected with the Holm step-down
procedure~\cite{holm1979simple}, which controls the family-wise error rate without assuming
independence between the tests.

Two decisions about how that family is read are taken here rather than after the results,
because either could otherwise be made to suit whatever the data turned out to be.

The first is which comparisons enter it, and the answer turns on the operating point rather than
the metric. Below saturation every variant completes work as fast as it is offered, so achieved
throughput there is pinned to the offered load by construction and cannot express a language
effect whatever the languages do. It is reported with its interval instead of tested. Above
saturation the reasoning inverts, because what each variant still sustains is pinned to nothing
and becomes the one quantity that can separate them, so a second operating point is fixed in the
overload region and sustained throughput and CPU utilisation are tested at it
(Section~\ref{sec:overload}). Two commitments go with that. The overload comparisons join the
same Holm family rather than forming a second one, so the division into two operating points
cannot be read as a split chosen once the $p$-values were known. And no latency percentile is
tested above saturation at all, because a producer that cannot keep up with a fixed intended
schedule falls behind cumulatively, which makes the percentiles there a lower bound and not an
estimate. Appendix~\ref{appendix:stats} sets out the reasoning behind both.

The second is what to conclude when the interval and the $p$-value disagree, which at modest
repeat counts they routinely will: a paired interval can sit wholly on one side of zero while
the corrected $p$-value stays far above any conventional threshold. Which instrument to believe
is not a free choice, and it does not depend on whether the metric happens to be interesting. It
depends on whether pairing removes the interference that measurement is exposed to. Pairing
cancels whatever an interleaved block shared, not interference that landed on one run of a pair
and not the other. A run-level summary such as peak resident set size is unmoved by a handful of
descheduled operations, and so is a central order statistic such as the median of tens of
thousands of per-operation samples, so there a paired difference that keeps its sign across
repeats is a difference in the software. p99 and p99.9 are the opposite case, computed from
precisely the operations that were disturbed, so an interval excluding zero is as consistent
with one run of a pair having been unlucky as with a real effect.
Table~\ref{tab:significance} is therefore blocked by that distinction, and intervals in its
second block are not read as findings however far from zero they sit. Applying the rule to the
block that suits the argument while quietly dropping it for the other is the specific failure
this pre-commitment exists to prevent.

The resolving power the design can reach is bounded by the number of repeats, and is stated
rather than left implicit. With \NRepeats{} repeats and \NTests{} comparisons the exact test's
two-sided floor is $\MinAttainableP$, and $\CorrectedFloor$ after correction, so a null reported
here reflects the data and not the arithmetic. Quoting both is what lets a reader tell a null
result from an unresolvable one. Choosing the repeat count so that the corrected floor sits below
the conventional $5\%$ is therefore part of the measurement protocol, alongside pinning the clock
and running the pinned cores tickless (Section~\ref{sec:testbed}).
Appendix~\ref{appendix:stats} derives the two floors.

\section{Reproducibility and research artifacts}\label{sec:repro}
The study is designed so that its results can be regenerated by a third party without access to
Nokia's proprietary code. Three things make that possible. The implementations are driven
entirely by the public frozen specification. The benchmarking harness, the fault catalogue and
its runner, and the FFI prototype are published with the thesis. And every result table, every
plot of measured data, and every quoted number in the evaluation is rebuilt by one regeneration
script in the published package (\texttt{bench/refresh\_thesis.sh}) directly from the raw logs.
Two things that claim might be read to cover are outside it. The hand-drawn schematics that
explain a concept rather than report a result are authored and not generated. And running the
same scripts over another host's campaign regenerates the artifacts by the same procedure, not
with the same values, which Section~\ref{sec:baremetal} demonstrates on a campaign of its own
design. What the script guarantees is that no figure, table, or sentence can hold a value the
logs no longer support.

The logs that script reads are part of the package. A performance campaign is an append-only
newline-delimited JavaScript Object Notation (JSON) file, named by the timestamp at which it
started. The reported sweep and the virtualised campaign it replaced both sit under
\texttt{bench/results/}, while the second bare-metal campaign has a directory of its own. Each
campaign has its own script and its own macro prefix, so a figure from one cannot reach a
sentence about another. The two bare-metal campaigns also write a sidecar holding every run's
complete histogram buckets, which is what the far-tail percentile is computed from and which the
virtualised campaign predates. Each run's environment is recorded by
\texttt{bench/capture\_env.sh} and stored beside the campaign data, and three runners write three
machine-readable outcome files under \texttt{faults/results/}: \texttt{taxonomy.csv},
\texttt{model\_checking.csv}, and \texttt{miri.csv}. Nothing in the evaluation chapters is
transcribed from those files by hand. The execution environment travels with its data, which
is what lets the thesis separate environment-sensitive results from
environment-insensitive ones. Nothing on the path requires proprietary source.

